\subsection{Kernergebnisse und Ausblick}\label{sec:kernergebnisse_und_ausblick}

Dieser Projektbericht hat mit Resteria ein Plattformkonzept für nachhaltigkeitsorientierte Hobbyköch:innen entwickelt, das eine konkrete Lücke im bestehenden Angebot schließt. Etablierte Rezept-Communitys bieten Austausch und Inspiration ohne dedizierten Zero-Waste-Fokus, Reste-Matching-Apps lösen umgekehrt das Zutaten-Problem ohne Ort für Austausch, wie~\autoref{sec:wettbewerbsanalyse} gezeigt hat; selbst leftovercooking bleibt bei der Community-Komponente auf das Verfolgen von Creator-Inhalten beschränkt. Resteria verbindet beide Seiten: Reste-Matching als Kernfunktion, ergänzt um Rettungspunkte, Reste-Level und einen Rettungs-Feed mit Challenges. Diese Differenzierung beruht auf der eigenen Wettbewerbsanalyse.

Die Gamification- und Community-Funktionen stützen sich dagegen auf literaturbasierte Verhaltenshebel: den Rahmen aus Drivers und Levers, der konkrete Funktionen statt bloßer Absichtserklärungen verlangt \parencite[S. 1]{vittuariHowReduceConsumer2023}, Reminders als wirksamen Nudge-Typ hinter der geplanten Erinnerungsbenachrichtigung \parencite[S. 10]{barkerWhatNudgeTechniques2021}, die belegte Wirksamkeit von Gamification gegen Lebensmittelverschwendung \parencite[S. 9]{somaFoodWasteReduction2020} sowie die Selbstwahrnehmungs-Mechanik hinter den Reste-Level-Titeln \parencite{nivedhithaHowGreenSocial2024}. Diese Trennung hält den Bericht davon ab, den gesamten Marktvorteil von Resteria als wissenschaftlich bewiesen darzustellen, obwohl er auf eigener Marktbeobachtung beruht.

Über den Rahmen dieses Projektberichts hinaus benötigt Resteria vor allem einen realen Nutzertest; die Phasenplanung in~\autoref{sec:phasenplanung} sieht ihn bereits als Beta-Test vor dem Launch vor. Daneben könnten die in~\autoref{sec:konzept} angelegten Kooperationen über ihren Erlösbeitrag hinaus neue Nutzer:innen erschließen, ohne den Fokus auf Reste-Matching und Community zu verwässern. Beide Schritte setzen voraus, dass das Konzept zunächst als Prototyp umgesetzt wird.
